\documentclass[10pt,a4paper]{amsart}
\usepackage[margin=19mm]{geometry}
\usepackage{amsmath,amssymb,mathtools}
\usepackage[scr=rsfs,cal=euler]{mathalfa}
\usepackage{xfrac}
\usepackage[unicode,hidelinks,bookmarks=false]{hyperref}
\newcommand{\ie}{i.e.\ }
\newcommand{\cf}{cf.\ }
\newcommand{\resp}{respectively\ }
\newcommand{\Subsection}{\subsection}
\providecommand{\nobreakdash}{\nobreak\hbox{-}}
\setlength{\emergencystretch}{2em}
\multlinegap0pt
\usepackage{bm}
\usepackage{bbm}
\usepackage{dsfont}
\usepackage{xspace}
\usepackage{cancel}
\usepackage{mathtools}
\usepackage{amsthm}
\makeatletter
\@ifundefined{lemma}{\newtheorem{lemma}{Lemma}}{}
\makeatother
\newcommand \R {{ \mathbb R}}
\title[Finite codimension stability of invariant surfaces]{Finite codimension stability of invariant surfaces}
\author{Giovanni Forni}
\date{Source: arXiv:2502.00898v2 (CC BY 4.0)}
\begin{document}
\maketitle

\noindent\emph{Source and licence.} Finite codimension stability of invariant surfaces. Version arXiv:2502.00898v2 (CC BY 4.0), distributed under the Creative Commons Attribution 4.0 International licence, as recorded in the arXiv licence metadata for this exact version and shown on the abstract page https://arxiv.org/abs/2502.00898v2. Full rights notice: licenses/arxiv-2502.00898-CC-BY-4.0.txt.\par\smallskip

\noindent\textit{Editorial scope.} Counted unit: the Lemma whose source label is \texttt{lemma:solutions} in the article (source0.tex lines 646--666, source label \texttt{lemma:solutions}), delivered together with the article's own complete original proof of it (source0.tex lines 667--728, one \texttt{proof} environment of 62 lines). No mathematics is written, repaired or re-derived: statement and proof are the article's own text line for line. Required context retained uncounted: none. Every definition, display and label of those lines is retained unchanged. Omitted from this package and not redistributed: 1--645, 729--909. Nothing inside a retained excerpt is omitted or altered: every retained body is delivered exactly as the article's lines are written, so no omission receipt of any kind accompanies this package. Citations: the retained excerpts carry no citation command at all, so no bibliography entry of the article is transcribed and no reference is redistributed. Reference adaptation: none was needed and none was made; every reference inside the delivered unit resolves inside it, so no unresolved reference or citation is produced. Printed slips kept literal: no printed word, symbol, hypothesis or number of the counted statement or of its proof was altered. Layout and class: the article's own class is \texttt{amsart} with packages including fontenc, lmodern, amsfonts, amsmath, amstext, amsbsy, amssymb, amsopn, amsthm, upref, eucal, mathptmx, mathrsfs, xcolor; the wrapper is the lane's pinned \texttt{amsart} editorial wrapper at 10pt on A4, which restates the theorem-like environments the retained text needs and adds the article's own macro definitions that the retained text needs (\texttt{\textbackslash R}), with extra wrapper packages (bm, bbm, dsfont, xspace, cancel, mathtools). The delivered numbering is the unit's own. Assets: no retained span contains a figure environment, a graphics-inclusion command, a PDF-page inclusion, a table or a TikZ picture; no image file is copied into this package, the package declares no asset dependency and no image file is redistributed with it. Licence: version arXiv:2502.00898v2 of this article is distributed under the Creative Commons Attribution 4.0 International licence, as recorded in the arXiv licence metadata for this exact version and shown on the abstract page https://arxiv.org/abs/2502.00898v2. A full rights notice is delivered at licenses/arxiv-2502.00898-CC-BY-4.0.txt. Characters: every retained excerpt is pure ASCII, so the delivered engine, XeLaTeX with its default Latin Modern text font, reports no missing character.
\begin{lemma}  
\label{lemma:solutions}
Let $\xi \in \R^2 \setminus \{0\}$ be such that the translation vector field $X_\xi$ on the translation surface $(M, \omega)$ is stable
(that is, the Lie derivative operator has close range) of finite codimension in Sobolev spaces of finitely differentiable functions, with loss of $\sigma>0$
derivatives in the scale $\{H^s_\omega(M) \}$ of weighted Sobolev spaces on $(M, \omega)$.
There are constants $\rho_1$, $\rho_2>0$ depending on $\Vert H \Vert_{C^3_\omega(M)}$  and a constant $K>0$ with the following property. If  $\Vert \mathcal F (H, u_0)\Vert_{C^1_\omega(M)} \leq \rho_1$, and the embedding $u: M \to M \times \R^2$  is such that $u\vert \Sigma = u_0 \vert \Sigma= \text{Id}_\Sigma$ and $\Vert u- u_0\Vert_{C^1_\omega(M)} \leq \rho_2$, then the linear para-homological equation in the unknown $(v, c)$,
$$
T_{M[u]} \begin{pmatrix} 0_2 & T_{S[u]} \\ 0_2 & 0_2 \end{pmatrix} T_{M[u]^{-1}}  v - 
T_{ M[u] } X_\xi \Big(T_{M[u]^{-1} }  v \Big)  + T_{M[u]}\Big( \sum_i  c_i [\chi_i] \Big) = f \,,
$$
has a linear solution operator 
$$
(v, c) = (v_1, v_2, c_1, c_2) =  (\mathcal L[u] (f),  P[u] (f)) =(\mathcal L_1[u] (f), \mathcal L_2[u] (f),  P_1[u] (f), P_2[u] (f) )
$$
such that the range of the operator  $\mathcal L[u]$ is contained in the subspace of functions vanishing at $\Sigma$, the range of the operator $P$
is finite dimensional, and the following estimate holds: for a function $\Phi$ increasing in all of its arguments, we have
$$
\Vert \mathcal L [u] (f) \Vert_{H^s_\omega(M)} + \vert P[u] (f) \vert \leq  C_s (K, \Vert H\Vert_{C^3_\omega(M)}, \xi) \Vert  f \Vert_{H^{s+2\sigma}_\omega(M)} \,.
$$
Moreover, the four linear operators of concern are all continuously differentiable mappings from $u\in C^1_\omega(M)$ to the space of linear operators (with operator norm). 
\end{lemma} 

\begin{proof}
We remark that by the definition of the matrix $S[u]$ we have
$$
\Vert S[u]  +I_2\Vert = \Vert S[u] - S[u_0] \Vert \leq C \Vert H \Vert_{C^3_\omega(M)} \Vert u-u_0 \Vert_{C^1_\omega(M)} \,.
$$
In fact, in our case we have that, 
$$
A[u_0] = \begin{pmatrix}   0_2   &   I_2 \\  0_2 & 0_2  \end{pmatrix} 
$$
and $Du_0^t =\begin{pmatrix} I_2& 0_2 \end{pmatrix}$, hence $L[u_0] =0$ and thus
$$
S[u_0] = \begin{pmatrix}  I_2 & 0_2 \end{pmatrix} \begin{pmatrix}   -I_2   &   0_2 \\  0_2 & I_2  \end{pmatrix}    \begin{pmatrix}  I_2 \\ 0_2 \end{pmatrix} = -I_2\,. 
$$
There exists $\rho_2>0$ such that for $\Vert u-u_0 \Vert_{C^1_\omega(M)} \leq \rho_2$, the para-products  $T_{M[u]}$ and $T_{M[u]^{-1}}$ are invertible with inverse 
such that 
$$  \Vert T^{-1}_{M[u]} -I \Vert_{H^s_\omega(M)} +  \Vert T^{-1}_{M[u]^{-1}} -I \Vert_{H^s_\omega(M)}   \leq 1/2\,.
$$
The equation can then be written as a system of two systems of 2 equations equations for the unknown vector-valued function 
$\hat v := T_{M[u]^{-1}}  v$:
$$
\begin{cases}  
T_{S[u]}  \hat v_2 -  X_\xi \hat v_1  + \sum_i  c_{i,1} [\chi_{i}]_1  =  (T^{-1}_{M[u]}  f )_1 \,, \\ \quad\quad \quad
-  X_\xi \hat v_2  + \sum_i  c_{i,2} [\chi_{i}]_2  =  (T^{-1}_{M[u]}  f )_2\,,
 \end{cases}
$$
This system can be solved by first solving the second equation, for an appropriate choice of the constants $\{c_{i,2}\}$. The solution
$\hat v_2$ which unique up to additive constants, can be plugged in the first equation, which can then be solved for an appropriate choice 
of the constants $\{c_{i,1}\}$.  More precisely the equation
$$
-  X_\xi \hat v_2  + \sum_i  c_{i,2} [\chi_{i}]_2  =  (T^{-1}_{M[u]}  f )_2
$$
can be solved in $H^{s+\sigma}_\omega(M)$  if (and only if) $c_{i,2} = D_i \Big( (T^{-1}_{M[u]}  f )_2\Big)$, for all invariant distributions 
$D_i$ in $H^{-(s+ 2\sigma)}_\omega(M)$, hence there exists a constant $C_{2,s}(\xi) >0$ such that 
$$
\sum_i \vert c_{i,2} \vert  \leq  C_{2,s} (\xi)  \Vert u \Vert_{C^1_\omega(M)} \Vert f \Vert_{H^{s+2\sigma}_\omega(M)}\,.
$$
The solution satisfies the estimate 
$$
\begin{aligned}
\Vert  v_2 \Vert_{H^{s + \sigma}_\omega(M)} &\leq  C'_{2,s} (\xi)  \Vert (T^{-1}_{M[u]}  f )_2- \sum_i  c_{i,2} [\chi_{i}]_2\Vert_{H^{s + 2\sigma}_\omega(M)}
\\ & \leq  C''_{2,s} (\xi)  \Vert u \Vert_{C^1_\omega(M)}  \Vert f \Vert_{H^{s + 2\sigma}_\omega(M)} \,.
\end{aligned}
$$
The first equation can be then solved in $H^s_\omega(M)$  if 
$$
c_{i,1} = D_i \Big (  (T^{-1}_{M[u]}  f )_1 - T_{S[u]}  \hat v_2\Big)\,,
$$
for all invariant distributions 
$D_i$ in $H^{-(s+ \sigma)}_\omega(M)$, hence there exists a constant $C_{1,s}(\xi) >0$ such that 
$$
\sum_i \vert c_{i,2} \vert \leq \Vert  (T^{-1}_{M[u]}  f )_1 - T_{S[u]}  \hat v_2   \Vert_{H^{s+\sigma}_\omega(M)}  \leq  C_{1,s} (\xi)  \Vert u \Vert_{C^1_\omega(M)} \Vert f \Vert_{H^{s+2\sigma}_\omega(M)}\,.
$$
The solution of the first equation then satisfies the bound
$$
\begin{aligned}
\Vert  v_1 \Vert_{H^{s}_\omega(M)} &\leq  C'_{2,s} (\xi)  \Vert  (T^{-1}_{M[u]}  f )_1 - T_{S[u]}  \hat v_2-  \sum_i  c_{i,1} [\chi_{i}]_1\Vert_{H^{s + \sigma}_\omega(M)}
\\ & \leq  C''_{2,s} (\xi)  \Vert u \Vert_{C^1_\omega(M)}  \Vert f \Vert_{H^{s + 2\sigma}_\omega(M)} \,.
\end{aligned}
$$
Finally, continuous differentiability of the operators for $u\in C^1_\omega(M)$  follows immediately from properties of para-products, that is, continuous differentiability of 
$T_a$ in $a$ and $T_{ab}- T_aT_ b$ with respect to $(a, b)\in (L^{\infty})^2$.
\end{proof}

\end{document}
